\documentclass[10pt,a4paper]{amsart}
\usepackage[margin=19mm]{geometry}
\usepackage{amsmath,amssymb,mathtools}
\usepackage[scr=rsfs,cal=euler]{mathalfa}
\usepackage{xfrac}
\usepackage[unicode,hidelinks,bookmarks=false]{hyperref}
\newcommand{\ie}{i.e.\ }
\newcommand{\cf}{cf.\ }
\newcommand{\resp}{respectively\ }
\newcommand{\Subsection}{\subsection}
\providecommand{\nobreakdash}{\nobreak\hbox{-}}
\setlength{\emergencystretch}{2em}
\multlinegap0pt
\usepackage{amssymb}
\newtheorem{lemma}{Lemma}
\newtheorem{proposition}{Proposition}
\title{On the mod $2$ cohomology algebra of oriented Grassmannians}
\author{Milica Jovanovi\'c, Branislav I.\ Prvulovi\'c}
\date{Source: arXiv:2306.11153v2 (CC BY 4.0)}
\begin{document}
\maketitle

\noindent\emph{Source and licence.} On the mod $2$ cohomology algebra of oriented Grassmannians. Version arXiv:2306.11153v2 (CC BY 4.0), distributed by arXiv under the Creative Commons Attribution 4.0 International licence, http://creativecommons.org/licenses/by/4.0/, as recorded in the arXiv licence metadata for this exact version and shown on the abstract page https://arxiv.org/abs/2306.11153v2. Full licence notice: \texttt{licenses/arxiv-2306.11153-CC-BY-4.0.txt}.\par\smallskip

\noindent\textit{Editorial scope.} Counted unit: the article's proposition kercapker2 (source0.tex lines 402--404). The article's complete original proof of it is delivered with it (source0.tex lines 405--462, one \texttt{proof} environment of 58 lines). No mathematics is written, repaired or re-derived: the statement and the proof are the article's own lines, in the article's own notation, and no retained line is substituted or re-flowed.  Retained uncounted as the source-local context the proof needs: lines 159--161; lines 332--341; lines 355--357; lines 359--361.  Nothing was omitted from any retained span, so the package carries no omission record.  No retained line contains a citation, so the wrapper carries no bibliography and no bibliography entry.  No retained line contains a figure environment, an image-inclusion command or drawing code, so no image is delivered and the package declares no asset dependency.  Editorial formatting only, in addition: an AMS \texttt{amsart} preamble that restates the article's own declarations and macros used by the retained lines, namely the environments \texttt{lemma}, \texttt{proposition} on separate counters, together with the standard packages those lines need.  The retained environments print in this document as Lemma 1, Proposition 1 in order of appearance, whereas the article prints the same items with the numbering of its own full document.  The complete native source of the article as distributed in the arXiv e-print for this exact version is bundled unchanged as \texttt{sources/143109/source0.tex}.
\begin{equation}\label{kohomologija_grasmanijana}
H^*(G_{n,k})\cong\mathbb Z_2[w_1,w_2,\ldots,w_k]/I_{n,k},
\end{equation}

\begin{lemma}\label{naturality}
{\rm(a)} For the map $i^*:H^*(G_{n+1,k})\rightarrow H^*(G_{n,k})$ we have
 \[i^*\big(w_r(\gamma_{n+1,k})\big)=w_r(\gamma_{n,k}), \quad 1\le r\le k.\]

 {\rm(b)} For the map $j^*:H^*(G_{n+1,k+1})\rightarrow H^*(G_{n,k})$ we have
 \begin{align*}
  j^*\big(w_r(\gamma_{n+1,k+1})\big)&=w_r(\gamma_{n,k}), \quad 1\le r\le k,\\
  j^*\big(w_{k+1}(\gamma_{n+1,k+1})\big)&=w_{k+1}(\gamma_{n,k})=0.   
 \end{align*}
\end{lemma}

\begin{equation}\label{w_rekurzivno}
    \overline w_{r+3} = w_1\overline w_{r+2} + w_2\overline w_{r+1} + w_3\overline w_r,\quad r\ge 0,
\end{equation}

\begin{equation}\label{overline_w}
    \overline w_r = \sum_{a+2b+3c = r} \binom{a+b+c}{a}\binom{b+c}{b} w_1^aw_2^bw_3^c,\quad r\ge 0,
\end{equation}

\begin{proposition}\label{kercapker2} Let $H^{2^t-1}(G_{2^t,3})\stackrel{w_1}\longrightarrow H^{2^t}(G_{2^t,3})$ be the multiplication with the Stiefel--Whitney class $w_1(\gamma_{2^t,3})$ and $i^*:H^{2^t-1}(G_{2^t,3})\rightarrow H^{2^t-1}(G_{2^t-1,3})$ the map induced by the embedding $i:G_{2^t-1,3}\hookrightarrow G_{2^t,3}$. Then
    $\ker w_1\cap \ker i^* = 0$.
\end{proposition}

\begin{proof}
    For a homogeneous polynomial $f\in\mathbb Z_2[w_1,w_2,w_3]$ we will denote by $[f]$ the cohomology class corresponding to $f$ via the isomorphism (\ref{kohomologija_grasmanijana}). If $[f]\in\ker w_1\cap \ker i^*$, then by Lemma \ref{naturality}(a) we first have $[f] = i^*[f] = 0$ in $H^{2^t-1}(G_{2^t-1,3})$, so 
    \[f = (\alpha w_1^2 + \beta w_2)\overline w_{2^t-3} + \gamma w_1\overline w_{2^t-2} + \delta \overline w_{2^t-1},\]
    for some $\alpha,\beta,\gamma,\delta\in\mathbb Z_2$. Since $\overline w_{2^t-2}, \overline w_{2^t-1}\in J_{2^t,3}$,
    \[[f] = [(\alpha w_1^2 + \beta w_2)\overline w_{2^t-3} + \gamma w_1\overline w_{2^t-2} + \delta \overline w_{2^t-1}] = [(\alpha w_1^2 + \beta w_2)\overline w_{2^t-3}],\]
    so we can choose $\gamma = \delta = 0$. 

    On the other hand, since $[f]\in \ker w_1$, we have $[w_1f] = 0$ in $H^{2^t}(G_{2^t,3})$, so
    \begin{equation}\label{w1f}
        w_1f = (\lambda w_1^2 + \mu w_2) \overline w_{2^t-2} + \nu w_1 \overline w_{2^t-1} + \varepsilon \overline w_{2^t},
    \end{equation}
    for some $\lambda, \mu, \nu, \varepsilon\in\mathbb Z_2$. By (\ref{w_rekurzivno}),
    \[\overline w_{2^t} = w_1\overline w_{2^t-1} + w_2\overline w_{2^t-2} + w_3\overline w_{2^t-3},\]
    and after substituting $(\alpha w_1^2 + \beta w_2)\overline w_{2^t-3}$ for $f$, (\ref{w1f}) becomes
    \begin{equation} \label{w1f1}
        (\alpha w_1^3 + \beta w_1 w_2 + \varepsilon w_3)\overline w_{2^t-3} = (\lambda w_1^2 + (\mu+ \varepsilon) w_2) \overline w_{2^t-2} + (\nu + \varepsilon)w_1\overline w_{2^{t-1}}.
    \end{equation}
By (\ref{overline_w}), the polynomial $\overline w_{2^t-2}$ contains the monomial $w_2^{2^{t-1}-1}$ with nonzero coefficient, and it is obvious that $(\mu+ \varepsilon) w_2\overline w_{2^t-2}$ is the only summand in (\ref{w1f1}) that can contain $w_2^{2^{t-1}}$. So we conclude $\mu+ \varepsilon=0$, i.e., $\mu = \varepsilon$. Equation (\ref{w1f1}) now becomes
    \begin{equation}\label{w1f2}
        (\alpha w_1^3 + \beta w_1 w_2 + \mu w_3)\overline w_{2^t-3} = \lambda w_1^2 \overline w_{2^t-2} + (\nu + \mu)w_1\overline w_{2^{t-1}}.
    \end{equation}
We now use the same idea to show that $\alpha =\beta = 0$. For $\alpha$ we observe the monomial $w_1^4w_2^{2^{t-1}-2}$, and use (\ref{overline_w}) in order to determine the coefficients of this monomial in each of the summands in (\ref{w1f2}). These coefficients are listed in Table \ref{tab:alpha}.
\begin{table}[h]
\renewcommand{\arraystretch}{1.5} %
\begin{tabular}{|c|c|c|c|c|}
 \hline
summand & $a$ & $b$ & $c$ &  $\binom{a+b+c}{a}\binom{b+c}{b}$ \\ [1ex]
 \hline
$\alpha w_1^3\overline w_{2^t-3}$ & 1 & $2^{t-1}-2$ & 0 &   1\\ 
$\beta w_1w_2\overline w_{2^t-3}$ & 3 & $2^{t-1}-3$ & 0 &   0\\
$\mu w_3\overline w_{2^t-3}$ & - & - & - &   -\\
$\lambda w_1^2 \overline w_{2^t-2}$ & 2 & $2^{t-1}-2$ & 0 &   0\\
$(\nu + \mu)w_1\overline w_{2^{t-1}}$ & 3 & $2^{t-1}-2$ & 0 &   0\\
 \hline
  \end{tabular}
\caption{\label{tab:alpha} Coefficients of $w_1^4w_2^{2^{t-1}-2}$ in (\ref{w1f2})}
\end{table}
So $\alpha w_1^3\overline w_{2^t-3}$ is the only summand containing $w_1^4w_2^{2^{t-1}-2}$, and its coefficient is $\alpha$. We conclude that $\alpha =0$.

For $\beta$ we proceed in the same way. We now observe the monomial $w_1^3w_2^{2^{t-1}-3}w_3$, and list its coefficients in Table \ref{tab:beta}, which allows us to conclude that $\beta =0$.
\begin{table}[h]
\renewcommand{\arraystretch}{1.5} %
\begin{tabular}{|c|c|c|c|c|}
 \hline
summand & $a$ & $b$ & $c$ &  $\binom{a+b+c}{a}\binom{b+c}{b}$ \\ [1ex]
 \hline
$\alpha w_1^3\overline w_{2^t-3}$ & 0 & $2^{t-1}-3$ & 1 &   0\\ 
$\beta w_1w_2\overline w_{2^t-3}$ & 2 & $2^{t-1}-4$ & 1 &   1\\
$\mu w_3\overline w_{2^t-3}$ & 3 & $2^{t-1}-3$ & 0 &   0\\
$\lambda w_1^2 \overline w_{2^t-2}$ & 1 & $2^{t-1}-3$ & 1 &   0\\
$(\nu + \mu)w_1\overline w_{2^{t-1}}$ & 2 & $2^{t-1}-3$ & 1 &   0\\
 \hline
  \end{tabular}
\caption{\label{tab:beta} Coefficients of $w_1^3w_2^{2^{t-1}-3}w_3$ in (\ref{w1f2})}
\end{table}

 Finally, we have $[f] = [(\alpha w_1^2 + \beta w_2)\overline w_{2^t-3}] = 0$ in $H^{2^t-1}(G_{2^t,3})$, and since $[f]$ was an arbitrary element of $\ker w_1\cap \ker i^*$, this concludes the proof.
\end{proof}

\end{document}
